\section{Progress and research-path record}

\subsection{Repository reset and traceability}

On 3 November 2026 the user requested a new research phase and report. The complete
former main tree is preserved under \texttt{archive/2026-11-03-pre-reset/}, with a named
archive branch at the original commit. The manifest also records the heads of 46
pre-existing branches; their unmerged contents remain separate and traceable.
Historical scope documents are not forward instructions for this phase.

The active tree records 21 idea IDs, seven families, experiments E00--E09, a claim ledger,
dated logbook, decision records and a CA2 evidence map. These distinguish what was known
before an experiment from what was learned after it. No narrative is reconstructed as
though an untested later idea motivated an earlier measurement.

\subsection{Completed engineering verification}

The reset's 10 unit tests passed locally. They check exact permutation and partial-sum
equivalence, signed reconstruction including the negative sign bit, complement identities,
count-support cardinality, deterministic comparator behaviour, split/metadata rejection
and non-overwriting records. The explicit E00 result and test transcript are curated
under \texttt{experiments/results/}.

\begin{table}[htbp]\centering\small
\begin{tabularx}{\linewidth}{@{}p{0.43\linewidth}X@{}}\toprule
Evidence & Current status / permissible interpretation\\\midrule
E00 arithmetic and metadata tests & Passed software checks; no physical or task-level conclusion.\\
E01 real workload screen & Exploratory independent digits256 screen recorded; frozen FYP model traces remain pending.\\
E02 target circuit/model & Contract/preflight and protocol prepared; physical inputs and validation missing.\\
E03--E08 target intervention comparisons & E03Q constructed ADC sensitivity run; no target physical intervention performed.\\
E09 Ising alternative & Held as a separate optional application, prerequisites unresolved.\\
Novelty / useful hardware improvement & Not established for any candidate.\\\bottomrule
\end{tabularx}
\caption{Status at the reset. Planned work is not reported as experimental success.}
\end{table}

\subsection{E03M: a constructed covariance check, 4 November}

E03M is a normalized probability model with \target{} binary storage bookkeeping,
four prospective 64-row groups and two input bit planes. It is not a target circuit
simulation, physical fan-in validation, real-workload demonstration or trained mapper.
The fixture uses declared current means 1 and 2 and independent 3\% Gaussian current
variation frozen across reads. No retired trial or foundry parameter is used.

Both mappings execute eight vector reads with 16 active rows and the same ideal count
in every column. Each receives the true full-group current deviation for scalar
calibration; midpoint count thresholds are ideal. Different shared-cell overlap gives
analog correlation patterns $[1,-1/3,1,-1/3]$ and $[1/3,1/3,1/3,1/3]$ after Gaussian
conditional-mean correction.

For decoded-error variance $v$ and covariance $c(\rho)$, the resulting MSEs are
\begin{equation}
 R_A=28v-8c(1/3),\qquad R_B=20v+16c(1/3).
\end{equation}
Their difference is positive for the constructed non-linear quantizer. With continuous
readout, $c(\rho)=v\rho$ and it vanishes. The mathematical expansion and proof, code,
configuration and original result hashes are retained in the research record.

\begin{table}[htbp]\centering\small
\begin{tabular}{@{}lrr@{}}\toprule
Constructed-model quantity & Mapping A & Mapping B\\\midrule
Executed vector reads & 8 & 8\\
Marginal count-error probability (\%) & 0.234308 & 0.234308\\
Analytic reconstructed MSE & 0.065222 & 0.047631\\
Continuous-readout MSE & 0.684 & 0.684\\\bottomrule
\end{tabular}
\caption{E03M, normalized arithmetic units squared. The 26.97\% analytic MSE
difference is an illustrative-model result, not a target-array improvement.}
\end{table}

Analytic identities were checked across all 256 columns. Paired Monte Carlo on one
predeclared column used three predeclared seeds with 131,072 independent variation
assignments each. Conditional-mean correction produced MSE reductions of 24.47--27.88\%;
separate full-group gain normalization retained 24.95--27.99\%. All paired difference
intervals exclude zero using per-assignment standard errors. This single constructed
setting establishes no advantage over a tuned CAR/adaptive-block frontier or arbitrary
pattern-specific calibration, and supplies no inference or energy result.

The result motivates a bounded prior-art and target pair-covariance check under ADR-0002.
It does not complete E03 or change the open final-topic decision. The curated record is
\path{experiments/results/E03M_20261104_covariance.json}; the method and proof are in
\path{research/novelty/2026-11-04-covariance-mapping.md}.

\subsection{E03Q: finite-ADC sensitivity (constructed-model pilot)}

A new E03Q software experiment extends the separate mathematical E03M fixture without
using any retired small-array measurements or foundry parameters. A \(256\times256\)
binary-weight \emph{bookkeeping} fixture retains four prospective 64-row groups,
two input bit cycles and one predeclared output column. Each read has 16 selected
rows with eight MTJs of each nominal current state. Manufacturing variation is
independent Gaussian with assumed 3\% cell-current coefficient of variation and
persists across both cycles. The same 131,072 independent evaluation-chip assignments
are paired across mappings; 65,536 separate chips select ADC-offset settings.
A shared, oracle-like full-group mismatch correction is inherited from E03M
and is not asserted physically feasible.

The finite ADC is a declared uniform \emph{midrise} quantizer on [16,32] normalized
current units followed by nearest quantized nominal references, lower-count tie-break
and clipping. Four through seven bits and phase offsets of
$-0.25$, 0 and $+0.25$ LSB are examined. The offset selected on calibration chips
is frozen before the independent evaluation.

\begin{table}[htbp]\centering\small
\begin{tabular}{@{}lrrr@{}}\toprule
Constructed decoder & Mapping A MSE & Mapping B MSE & Relative decrease (\%)\\\midrule
Ideal midpoint & 0.068687 & 0.048462 & 29.4\\
Fixed 6-bit ADC & 2.261963 & 1.917000 & 15.3\\
Tuned 6-bit ADC & 0.122276 & 0.087631 & 28.3\\
Tuned 7-bit ADC & 0.081520 & 0.057892 & 29.0\\\bottomrule
\end{tabular}
\caption{E03Q held-out, one-column Monte Carlo. Normalized integer-output MSE;
no circuit, inference or energy result.}
\end{table}

The paired difference (A minus B) for the tuned 6-bit setting is 0.03465,
with an approximate chip-level 95\% interval [0.02906,0.04023] under independent
chip sampling. The fixed 4- and 5-bit operating points exhibit strong
threshold-alignment artefacts, illustrating how calibration can dominate the
comparison. The mappings have identical marginal Gaussian current laws by
construction; their finite-sample empirical read-error frequencies are not exactly
identical. This is a useful decoder-sensitivity check, but not a new MRAM mechanism
or evidence of a useful trained mapper. All model inputs and output hashes are
registered with the E03Q pilot.

\subsection{External task-level leads and reproduction boundary}

The latest externally supplied draft mentions neural mapping and noisy-Ising
results from another discussion. Their raw outputs and source implementations were
not independently reproduced here. They remain unverified external leads; no
claimed agreement in outcome establishes a shared mechanism. Likewise, an
incremental cached-field Ising error is not the same operator as full analog
field recomputation.

\subsection{E09A follow-up: exact Ising controls}

The revision-15 consolidation added six meaningful software tests; the combined suite
has 16 passing tests. E09A verifies small-graph Max-Cut optima, the target-energy/cut
identity, full versus incremental sequential heat-bath trajectories under common random
inputs, and exact cache parity after every update. Eight graph/seed cases include a
64-spin arithmetic-only fixture. No noisy-target success-rate benchmark was executed.

For one coupling coefficient, flipping a spin and reversing it leaves an incremental
read-error residual $2s_k^{\rm old}(\eta_2-\eta_1)$. Identical frozen coefficient errors
cancel; independent zero-mean errors of variance $\sigma^2$ give residual variance
$8\sigma^2$. An exact unit-sign enumeration checks this algebra. It is a cache-control
diagnostic, not a measured noise law, new general theorem or target solver result.
The target-kernel preflight remains blocked by missing response and temporal-noise
evidence, and the scientific adapter is not implemented. Curated evidence is under
\path{experiments/results/E09A_20261104_reference.json}.

\subsection{E01R: resource audit and executable NUS handoff}

The archived neural software identifies a reusable 784--128--64--10 MNIST model,
frozen W4A8 quantizer, saved split indices and raw dataset locations. Their current
binary presence and hashes require checking on NUS. The new runner automates those
checks and verifies its export against the original, source-pinned PyTorch quantizer.
It does not retrain, alter calibration or import retired electrical results.

The first protocol uses 128 saved calibration parents and 256 previously inspected
validation parents, explicitly labelled development. All four FC1 row blocks and all
four weight planes are exported. Each block packs 128 outputs and four weight bits
into two prospective \target{} binary position maps; eight arrays are needed across
the layer. The final block uses 16 rows and has 240 padding rows. These are capacity
and address maps, not physical validation. Parent-image/activation-bit identities,
signed reconstruction and one common resident mapping across the planes are retained.

Twenty-seven software tests pass, including eleven new export/resource checks and
an integration fixture whose external NUS loader is stubbed. The local resource CLI
correctly reports the unavailable NUS checkpoint/data and exits 2. No actual frozen-
MNIST golden-forward run, primary E01 workload result or fresh target circuit run has
been completed here. The ready command and return-packet protocol are in
\path{START_HERE.md}; E01 remains pending user execution in the authorized environment.

\subsection{How results enter CA2}

Each meaningful run records its hypothesis, code revision, command, source hashes,
calibration/evaluation roles, chip identities, evidence class, outputs, uncertainty and
decision. The experiment register tracks readiness and blocking evidence. A dated
decision updates the claim ledger and relevant report section. Raw runs are never
overwritten; curated evidence links the analysis to the original outputs.
